% Candidate prose only. Not included in the active manuscript.
% Figure/table insertion and numbering await author review.
\section{Raw-Gradient Information Along Adam Trajectories}

Adam transforms gradients through first-moment averaging and coordinatewise
scaling before updating the parameters. Raw-gradient reversal therefore does
not directly describe reversal of the applied updates. Nevertheless, successive
raw gradients measure how the training objective changes along the trajectory
chosen by Adam. We ask whether this response contains information about an
upcoming loss excursion beyond the loss history already observed.

We retain the projected-change expression
\begin{equation}
  \eta\widehat\kappa_t
  =1-\rho_t\frac{\|g_t\|}{\|g_{t-1}\|}.
\end{equation}
Under Adam, this is a normalized projected gradient change, without an
identified Hessian-curvature interpretation or stability threshold. For nonzero
gradients, a value above two requires both negative consecutive cosine and
gradient-norm growth. Showing these quantities together explains the observed
geometry; it does not provide independent corroboration of three mechanisms.
The stored first moment supplies a complementary view of how Adam averages
the evolving gradient input.

\subsection{Controlled Trajectories and Local Excursions}

We examine CIFAR-10/MLP, CIFAR-10/ResNet-32 and WikiText-2/Transformer under
full-batch Adam, using one frozen initialization per case and three constant
learning rates. Each trajectory contains 1,024 updates, with fresh Adam
moments, $(\beta_1,\beta_2)=(0.9,0.999)$ and $\epsilon=10^{-8}$.
There is no clipping, weight decay, schedule or learning-rate intervention.
The overview displays loss, projected change, raw-gradient cosine and
gradient/momentum norms. Its shorter trailing average is a display reduction;
all timing and prediction analyses use the unsmoothed observations.

The micro views show each highest-rate arm's completed episode with the
largest loss peak relative to its episode baseline. These explicitly selected
examples illustrate the temporal sequence. At the Transformer peak at state 794,
projected change first crosses two within the local pre-onset window at 788,
when loss is only 0.40\% above its recent minimum. Loss reaches 12.0\% above that
minimum at 790. Gradient and momentum norms also grow before this pronounced
increase. Projected change, gradient norm, momentum norm and loss attain their
local maxima at 790,792,793 and 794, respectively. The crossing precedes the
pronounced rise, but not its first small positive loss change. This exact
ordering of maxima is not shared by every case.

\subsection{Recurrence, Lead Time and Specificity}

To assess repetition, we define episodes using loss alone. An episode begins
at a 10\% increase above the preceding 16-state minimum and ends at the first
return below that frozen baseline. Rebounds before recovery remain one
episode; unfinished recoveries are retained. The nine arms contain 129
episodes,45 of which reach at least 50\% above baseline.

The episode table reports every arm, including misses and zero-event arms.
At the highest rates, a fresh crossing of two occurs 1--8 updates before the
10\% onset, while loss is still within 5\% of its recent minimum, in 13/17 large
MLP episodes,3/3 ResNet episodes and 7/7 Transformer episodes. Among episodes
with a warning, median leads are 5,2 and 5 updates, respectively. Coverage is
lower when all episode sizes are considered:17/39,5/7 and 13/13. The
onset-aligned traces show the individual trajectories behind these counts,
rather than only a single selected peak or a mean trajectory.

A recurring precursor need not be a specific warning. On the forward-test
states, projected change exceeds two at 891 eligible states, of which only 187
are followed within 8 updates by a loss at least 10\% above the current value.
Thus its precision is 21.0\%; a fresh crossing has precision 15.9\%.
The smallest-rate ResNet and Transformer arms have no qualifying episode,
yet contribute 175 and 190 above-two alarm states with no positive forecast
targets. These controls preclude treating two as a universal Adam instability
boundary. Episode coverage and alarm precision use different denominators
and different loss references, which are specified in the methods.

\subsection{Information Beyond Recent Loss}

We next test incremental information, using the same eligible states for all
models. The target is a 10\% loss increase within the next 8 updates, from a
state still within 5\% of its trailing minimum. The baseline includes task/rate
indicators, training age, current loss and recent loss changes. We compare it
with the same regularized logistic model plus continuous projected change.
Five expanding-prefix fits predict later time blocks, with a 24-state exclusion
gap; scaling uses training data only.

Across 4,534 forward-test decisions, including 495 positives, average precision
increases from 0.260 to 0.454. Brier score decreases from 0.1197 to 0.0868.
Both improve in all five time blocks, and the per-case comparisons retain the
direction of improvement. These case summaries subset pooled-model predictions;
they do not measure transfer to an unseen task. The low ResNet event prevalence
and the numbers of decisions and positives are reported with the scores.

Additional checks clarify the finding. Raw cosine alone yields average
precision 0.423. Expanding the loss baseline to all 16 recent log-loss changes
retains an improvement from 0.258 to 0.447 after adding projected change, although
AP improves in four rather than all five blocks. With a common larger exclusion
gap, current projected change is more useful in aggregate than its 16- or 32-state
delayed values. The complete ablations and target/horizon sensitivities are
retained; adding all diagnostics is not uniformly better.

These comparisons establish incremental predictive information under the
tested models, rather than calibrated risk estimates. The original loss-only
model has worse pooled Brier score than a training-prevalence predictor
(0.1197 versus 0.0986). The projected extension improves the pooled comparison,
but does not beat that probability reference within every case or gap setting.
The evidence consequently supports a diagnostic association along the measured
trajectories, with no validated operating threshold or scheduling benefit.

The sequence is scientifically useful because raw-gradient geometry remains
informative even when Adam transforms those gradients before applying an
update. The local views describe the response, the episode census establishes
its recurrence and failures, and the forward comparison tests information
beyond recent loss. The analysis remains exploratory, with dependent events
from one initialization per case. It neither identifies a causal feedback
mechanism nor establishes an intervention effect.
